\section{Conclusions}

In this work we have developed a tensor-based framework for the
analysis and transformation of centrifugal distortion parameters in
high-resolution rotational spectroscopy.  The central idea is to
regard the reduced quartic and sextic tensors as the primary objects,
and to interpret the Watson centrifugal distortion constants as
coordinates obtained by projection onto a particular effective
Hamiltonian parametrization.

The tensorial viewpoint itself is not new: it is already implicit in
classical rovibrational perturbation theory and in the effective
Hamiltonian formalism of Watson and related work.  The novelty of the
present formulation lies instead in making this structure explicit as
the basis of a practical framework that connects the direct and inverse
problems.  On the direct side, tensor quantities can be obtained from
quantum-mechanical force fields and rovibrational couplings.  On the
inverse side, the same tensor quantities can be reconstructed from
spectroscopic Watson constants by pseudoinverse reconstruction.

For the quartic case, the natural reduced tensor is the symmetric
$3\times 3$ matrix $\Tau'$, obtained from the pair--pair quartic
tensor.  This tensor admits a spectral decomposition in terms of three
eigenvalues and three orientation coordinates, thereby providing a
clear geometric description of quartic centrifugal distortion prior to
the final choice of Watson parametrization.  The passage from the
tensorial description to the five spectroscopic quartic constants is
then understood as a projection from tensor space to the effective
Hamiltonian coordinates.  Conversely, the inverse reconstruction from
Watson constants is an affine gauge problem, in which the
Moore--Penrose pseudoinverse selects one distinguished representative
within the one-parameter family of equivalent quartic tensors.

This interpretation leads naturally to a tensor-based algorithm for
representation transforms.  Starting from Watson centrifugal
distortion constants obtained either from spectroscopic fits or from
quantum-chemical calculations, one reconstructs the reduced tensor,
permutes the molecular axes, and finally reprojects the transformed
tensor onto Watson constants in the new representation.  For quartic
centrifugal distortion this procedure reproduces the analytical
transformation formulas of Yamada and Winnewisser with machine
precision.  For sextic centrifugal distortion, where analogous
representation transforms are not generally available in closed form,
the same reconstruction strategy provides a practical and numerically
stable route to internally consistent transformations.

The usefulness of the method was illustrated for dimethylsulfoxide and
for a series of HBN isomers.  In the DMSO case, the analysis confirms
that large changes in individual Watson constants under representation
transforms arise primarily from the parametrization of the effective
Hamiltonian and from the numerical conditioning of the chosen
reduction--representation combination.  In the HBN series, the tensor
framework provides a more robust basis for comparing the quartic
centrifugal distortion field across related isomers than the Watson
constants themselves, because the comparison can be carried out at the
level of tensor invariants rather than at the level of
representation-dependent fitted parameters.

An important practical consequence of the present framework is that
representation transforms can be applied directly to spectroscopic
parameter sets, without requiring explicit analytical transformation
matrices tailored to each reduction and axis system.  The same
approach also provides a natural interface between spectroscopic fits,
quantum-chemical calculations based on harmonic and cubic force
fields, and direct rovibrational calculations.  In this sense, the
tensor reconstruction procedure establishes a common tensorial language
for centrifugal distortion parameters obtained from heterogeneous
sources of information.

A macOS application implementing the quartic and sextic
representation-transform algorithms discussed in this work is
available from the author upon request.

The tensor formulation developed here should therefore be viewed not
simply as an alternative parametrization of Watson's Hamiltonian, but
as a framework that makes explicit the relation between tensorial
rovibrational content and spectroscopic effective-Hamiltonian
coordinates.  This provides both a clearer conceptual picture of
centrifugal distortion and a practical computational tool for
performing representation transforms, analyzing tensor invariants, and
comparing spectroscopic and quantum-chemical descriptions on the same
footing.
